\chapter{Conclusions}
\label{ch:conclusions}

% The project is a success. Summarise what you have done and accomplished.
结合log文件和 3.3 节，计算速度快又计算准确的是自适应步长的方法，其次是RK方法，再其次是LMM方法。对于固定步长的方法，隐式的方法比显式的要慢（程序中设定了不动点迭代超过一万次时，
若前后两项的误差还没有小于机器精度，则停止迭代），同一求解器高阶所用的时间变长，但是计算准确率大幅提高。

当误差很小时，从 3.4 max-norm 可以看出加大步长，误差减小并不符合精度，推测是此时算法的误差不占主要地位。
当误差很大时，绘制的折线也不符合，结合 3.5 节可知是因为算法并没有较为准确的绘制图像。

在启动LMM算法前，为了获取足够的数据，我一开始使用Euler方法，效果非常差，后来使用了ClassicalRK方法，得到明显改观，
这也体现出要求解的两个问题对初始条件非常敏感。

程序设计的不足之处在于，
\begin{itemize}
    \item 一开始没有认识到其实所有求解器的主要不同在于如何从上一个step的数据得到下一个step的数据，导致有许多的重复代码。
    \item 对于自适应步长的方法如何设定其$E_{abs},E_{rel},\rho,\rho_{max},\rho_{min}$没有好的想法，直接硬编码到其中，给与之相关的测试带来麻烦。
    \item 对于确定自适应步长算法确定满足误差小于0.001的最小cpu time的参数，没有设计好的算法。
    \item 对于控制自适应步长算法的总步数，没有设计好的算法，采用了直接遍历区间的方法，并且找到的参数绘制出来的图片效果不理想，可能存在更好的参数。
\end{itemize}
